\documentclass[a4paper]{article}
\usepackage{times}
\usepackage[utf8]{inputenc}
\usepackage[T1]{fontenc}
\usepackage{selinput}
\usepackage{upquote}
\usepackage[margin=2cm, rmargin=4cm, tmargin=3cm]{geometry}
\usepackage{tcolorbox}
\tcbuselibrary{breakable}
\usepackage{xspace}
\usepackage{amsmath}
\usepackage{float}
\usepackage[french]{babel}
\usepackage{hyperref}
\usepackage{xurl}
\usepackage{fontawesome6}
\usepackage{enumitem}
\usepackage{marginnote}
\usepackage{ulem}
\usepackage[yyyymmdd,hhmmss]{datetime}
\usepackage{ulem}
\usepackage{environ}
\usepackage{graphicx}
%\usepackage[top=Bcm, bottom=Hcm, outer=Ccm, inner=Acm, heightrounded, marginparwidth=Ecm, marginparsep=Dcm]{geometry}
\usepackage{courier}



\newtcolorbox{Example}[1]{colback=white,left=20pt,colframe=slideblue,fonttitle=\bfseries,title=#1}
\newtcolorbox{Solutions}[1]{breakable,colback=white,left=20pt,colframe=green,fonttitle=\bfseries,title=#1}
\newtcolorbox{Conseils}[1]{colback=white,left=20pt,colframe=slideblue,fonttitle=\bfseries,title=#1}
\newtcolorbox{Warning}[1]{colback=white,left=20pt,colframe=warning,fonttitle=\bfseries,title=#1}
\newtcolorbox{Note}[1]{colback=white,left=20pt,colframe=grey,fonttitle=\bfseries,title=#1}
\newtcolorbox{Dessin}[1]{colback=white,left=20pt,colframe=grey,fonttitle=\bfseries,title=#1}

\setlength\parindent{0pt}

  %Exercice environment
  \newcounter{exercice}
  \newenvironment{Exercice}[1][]
  {
  \par
  \stepcounter{exercice}\textbf{Question \arabic{exercice}:} (\faClock \enskip \textit{#1})
  }
  {\bigskip}
  

% Title
\newcommand{\titre}{\begin{center}
  \section*{Algorithmes et Pensée Computationnelle}
\end{center}}
\newcommand{\cours}[1]
{\begin{center} 
  \textit{#1}\\
\end{center}
  }



% \newcommand{\exemple}[1]{\newline~\textbf{Exemple :} #1}
\newcommand{\todo}[1]{\textcolor{red}{#1}}
%\newcommand{\attention}[1]{\newline\faExclamationTriangle~\textbf{Attention :} #1}

% Documentation url (escape \# in the TP document)
\newcommand{\documentation}[1]{\faBookOpen~Documentation : \href{#1}{#1}}

% Clef API
\newcommand{\apikey}[1]{\faKey~Clé API : \lstinline{#1}}
\newcommand{\apiendpoint}[1]{\faGlobe~Url de base de l'API \href{#1}{#1}}

%Listing Python style
\usepackage{color}
\definecolor{slideblue}{RGB}{33,131,189}
\definecolor{green}{RGB}{0,190,100}
\definecolor{blue}{RGB}{121,142,213}
\definecolor{grey}{RGB}{120,120,120}
\definecolor{warning}{RGB}{235,186,1}

\usepackage{listings}
\lstdefinelanguage{texte}{
    keywordstyle=\color{black},
    numbers=none,
    frame=none,
    literate=
           {é}{{\'e}}1
           {è}{{\`e}}1
           {ê}{{\^e}}1
           {à}{{\`a}}1
           {â}{{\^a}}1
           {ù}{{\`u}}1
           {ü}{{\"u}}1
           {î}{{\^i}}1
           {ï}{{\"i}}1
           {ë}{{\"e}}1
           {Ç}{{\,C}}1
           {ç}{{\,c}}1
           {ô}{{\^o}}1,
    columns=fullflexible,keepspaces,
	breaklines=true,
	breakatwhitespace=true,
}
\lstset{
  language=Python,
	%basicstyle=\bfseries\footnotesize,
  basicstyle=\footnotesize\ttfamily\small,
	breaklines=true,
	breakatwhitespace=true,
	commentstyle=\color{grey},
	stringstyle=\color{slideblue},
  keywordstyle=\color{slideblue},
	morekeywords={with, as, True, False, Fraction, Float, join, index, None, main, argparse, self, sort, __eq__, __add__, __mul__ ,__ne__, __radd__, __del__, __ge__, __gt__, split, os, endswith, is_file, scandir, @classmethod, this},
	deletekeywords={id},
	showspaces=false,
	showstringspaces=false,
	columns=fullflexible,keepspaces,
	literate=
           {é}{{\'e}}1
           {è}{{\`e}}1
           {ê}{{\^e}}1
           {à}{{\`a}}1
           {â}{{\^a}}1
           {ù}{{\`u}}1
           {ü}{{\"u}}1
           {î}{{\^i}}1
           {ï}{{\"i}}1
           {ë}{{\"e}}1
           {Ç}{{\,C}}1
           {ç}{{\,c}}1
           {ô}{{\^o}}1,
    numbers=left,
}

\newtcbox{\mybox}{nobeforeafter,colframe=white,colback=slideblue,boxrule=0.5pt,arc=1.5pt, boxsep=0pt,left=2pt,right=2pt,top=2pt,bottom=2pt,tcbox raise base}
\newcommand{\projet}{\mybox{\textcolor{white}{\small projet}}\xspace}
\newcommand{\optionnel}{\mybox{\textcolor{white}{\small Optionnel}}\xspace}
\newcommand{\advanced}{\mybox{\textcolor{white}{\small Pour aller plus loin}}\xspace}
\newcommand{\auto}{\mybox{\textcolor{white}{\small Auto-évaluation}}\xspace}


\usepackage{environ}
\newif\ifShowSolution
\NewEnviron{solution}{
  \ifShowSolution
	\begin{Solutions}{\faTerminal \enskip Solution}
		\BODY
	\end{Solutions}
  \fi}

\usepackage{environ}
\newif\ifShowExample
\NewEnviron{exemple}{
  \ifShowExample
	\begin{Example}{\faTerminal \enskip Exemple}
		\BODY
	\end{Example}
  \fi}

\usepackage{environ}
 \newif\ifShowConseil
  \NewEnviron{conseil}{
    \ifShowConseil
    \begin{Conseils}{\faLightbulb \quad Conseil}
      \BODY
    \end{Conseils}

    \fi}

    \usepackage{environ}
  \newif\ifShowWarning
  \NewEnviron{attention}{
    \ifShowWarning
    \begin{Warning}{\faTriangleExclamation \quad Attention}
      \BODY
    \end{Warning}

    \fi}

  \newif\ifShowNote
  \NewEnviron{note}[1]{
    \ifShowNote
    \begin{Note}{\faEye \quad Comment added by #1 at \date - \currenttime}
      \BODY
    \end{Note}

    \fi}  

    \newif\ifShowDessin
    \NewEnviron{dessin}[1]{
      \ifShowDessin
      \begin{Dessin}{\faPen \quad Figure}
        \BODY
      \end{Dessin}
  
      \fi}  

%\newcommand{\Conseil}[1]{\ifShowIndice\ \newline\faLightbulb[regular]~#1\fi}

\usepackage{array}
\usepackage[T1]{fontenc}

\newcolumntype{C}[1]{>{\centering\let\newline\\\arraybackslash\hspace{0pt}}m{#1}}
\setlist[itemize]{label=\textbullet}

\begin{document}
% Change the following values to true to show the solutions or/and the hints
\ShowSolutiontrue
\ShowConseiltrue
\titre
\cours{Algorithmes de recherche - Exercices avancés}

Le but de cette séance est de se familiariser avec les algorithmes de recherche. Dans cette série d'exercices, nous manipulerons des listes et collections en Java et Python. Nous reviendrons sur la notion de récursivité et découvrirons les arbres binaires de recherche. Au terme de cette séance, l'étudiant sera en mesure d'effectuer des recherches de façon efficiente sur un ensemble de données. \\

Le code présenté dans les énoncés se trouve sur Moodle, dans le dossier \lstinline{advanced}.


\section{Recherche binaire}

% ADVANCED
\begin{Exercice}[15 minutes] [Python] Recherche binaire - plus proche élément\\

Soit une liste d'entiers triée \lstinline{liste} et un entier \lstinline{e}. Écrivez un programme retournant la valeur dans \lstinline{liste} la plus proche de \lstinline{e} en utilisant une recherche binaire (binary search). On suppose que la liste n'est pas vide. \\

\lstinputlisting{code/advanced/plus_proche_binaire.py}

\begin{conseil}
    Pensez à définir des variables \lstinline{low} et \lstinline{high} délimitant l'intervalle de recherche. 
    
\end{conseil}

    \begin{solution}
        \lstinputlisting[firstline=1, firstnumber=1, lastline=20]{solutions/advanced/plus_proche_binaire.py}
    \end{solution}
    \begin{solution}
        \lstinputlisting[firstline=21, firstnumber=21]{solutions/advanced/plus_proche_binaire.py}
    \end{solution}

\end{Exercice}



% ADVANCED
\begin{Exercice}[20 minutes] [Python] Recherche dans une matrice %\optionnel\\
    
    Soit une matrice ordonnée \lstinline{matrice} de $n$ lignes et $m$ colonnes, et un entier \lstinline{e}. Une matrice est ordonnée si :
    \begin{itemize}
        \item Chaque ligne est triée par ordre croissant, de gauche à droite : \lstinline{matrice[i][j] <= matrice[i][j+1]} ;
        \item Chaque colonne est triée par ordre croissant, de haut en bas : \lstinline{matrice[i][j] <= matrice[i+1][j]}.
    \end{itemize}

    Écrivez une fonction \lstinline{recherche_matricielle(matrice, e)} qui retourne la position \lstinline{(ligne, colonne)} de \lstinline{e} dans \lstinline{matrice}. Si \lstinline{e} apparaît plusieurs fois, retournez l'une de ses positions, peu importe laquelle. Si \lstinline{e} n'est pas dans la matrice, retournez \lstinline{(-1, -1)}.\\

    \begin{Example}{\faTerminal \quad Exemples}
        \textbf{Exemple 1}\!\!\!\!: pour \lstinline{matrice = [[1,2,3,4],[4,5,7,8],[5,6,8,10],[6,7,9,11]]} et \lstinline{e = 7}, la fonction retourne \lstinline{(1, 2)} ou \lstinline{(3, 1)} : 7 apparaît deux fois, une seule position suffit.\\

        \textbf{Exemple 2}\!: pour \lstinline{matrice = [[1,2],[3,4]]} et \lstinline{e = 7}, la fonction retourne \lstinline{(-1, -1)} car 7 n'est pas dans la matrice.
    \end{Example}

    \lstinputlisting{code/advanced/recherche_matricielle.py}

    \begin{conseil}
    Commencez par une solution simple : deux boucles \lstinline{for} imbriquées (une boucle \lstinline{for} dans une autre) permettent de parcourir toutes les cases d'une matrice. Demandez-vous ensuite comment tirer parti du fait que chaque ligne est triée pour éviter de tout parcourir.
    \end{conseil}

    \begin{solution}
        \textbf{Version 1: parcours de toutes les cases}
        \lstinputlisting[firstline=1, firstnumber=1, lastline=8]{solutions/advanced/recherche_matricielle.py}

        \textbf{Complexité : $O(n \cdot m)$.} On examine les cases une par une. Dans le pire des cas (\lstinline{e} absent de la matrice ou situé dans la dernière case), on examine les $n \cdot m$ cases. Cette version n'utilise pas du tout le fait que la matrice est ordonnée : elle fonctionnerait aussi sur une matrice quelconque.
    \end{solution}
    \begin{solution}
        \textbf{Version 2: recherche binaire dans chaque ligne}
        \lstinputlisting[firstline=9, firstnumber=9]{solutions/advanced/recherche_matricielle.py}

        \textbf{Complexité: $O(n \log m)$.} Chaque ligne étant triée, on y cherche \lstinline{e} par recherche binaire : à chaque itération, l'intervalle \lstinline{[low, high]} est divisé par deux, donc une recherche dans une ligne de $m$ éléments demande au plus environ $\log_2 m$ itérations, soit $O(\log m)$. Dans le pire des cas, on répète cette recherche pour chacune des $n$ lignes, d'où $O(n \log m)$ au total.\\

        \textbf{Remarque:} on ne peut pas faire une seule recherche binaire sur toute la matrice (par exemple pour trouver la bonne ligne, puis chercher uniquement dans celle-ci), car les lignes peuvent se chevaucher : le dernier élément d'une ligne n'est pas forcément plus petit que le premier de la ligne suivante. Dans l'exemple 1, la ligne \lstinline{[4,5,7,8]} se termine par 8, alors que la ligne suivante commence par 5, et 8 apparaît dans ces deux lignes.
    \end{solution}

\end{Exercice}

\end{document}
